% ===========================================================================
% 05-Limitations.tex.
% ===========================================================================

\chapter{Limitations and future research directions}
\label{sec:limitations}
\framedecl{ER}

\begin{question}
Every measurement in this thesis was made on one atlas, one checkpoint and one
slice of both. Which of its conclusions are properties of the model, and which
are properties of that slice?
\end{question}

\noindent
A limitation is worth writing down when it changes what a number licenses.

% ===========================================================================
\section{Populations and selection effects}
\label{sec:lim-population}

\paragraph{Only reproducible conditions} A residual that does not reproduce
between two halves of the same well is not a target anything could learn, so the
residual arm trains on the $50.5\%$ of conditions whose residual does reproduce,
$\num{2523}$ of $\num{4999}$ training conditions kept at $\cos > 0.109$, the
$95$th percentile of a null pairing half A of one condition with half B of
another.\gloss{reliability line}{The half-split cosine threshold above which a
condition's residual counts as reproducible.} The model was taught the
empirically reproducible fraction. Split-half retrieval establishes
discriminability of treated populations under this sampling design and cell
budget; it does not establish predictability from untreated input or
learnability.

A reliability filter might be expected to select \emph{toward}
context-in\-de\-pen\-dent conditions and bias the transfer coefficient upward.
It does the opposite. Unfiltered \ci{0.597}{0.563}{0.632}, filtered
\ci{0.553}{0.514}{0.593}; the filtered value is quoted, because disattenuation
is unstable at the very low reliabilities the filter removes.

That calculation's own filter is the inherited $\cos > 0.2$ line, where the
training set uses $0.109$, so the two estimates bracket the trained-on
population without reproducing it. The generation evaluation mismatches
differently again, applying $0.109$ to \texttt{train} and leaving the held-out
splits unfiltered. That one is fundamental and is taken up below.

\paragraph{The splits are scored on different populations} The within-well
split-half precision reference differs sharply across the three splits as
scored, $0.956$ on \texttt{train}, $0.824$ on \texttt{unseen\_combo} and $0.836$
on \texttt{unseen\_drug}. The obvious reading is that held-out conditions carry
less recoverable signal. That reading is wrong.

It is the selection asymmetry noted above. \texttt{train} is drawn from a pool
already filtered at the $0.109$ line, and none of its $300$ scored conditions
falls at or below it, while the held-out splits are unfiltered and still contain
conditions whose two halves disagree outright. Impose the training line on all
three and the references agree at $0.956$, $0.961$ and $0.960$, a spread of
$0.005$ where there had been $0.132$ (\cref{fig:splitref}). Matched on
reproducibility, a held-out condition is as recoverable as a trained one.

% The float that stood here (193-249 of the previous revision) is now
% figs/splitref.tex. It moved out of the section for the same reason the
% other reworked figures did -- the drawing carries 70 lines of its own
% rationale, which does not belong in the prose -- and it lost its
% \begin{fullwidth} on the way: it reserved 174mm, drew 150.2mm, and cost
% this page its margin column to show six numbers the caption already states.
% It now draws 132.9mm inside the 142mm measure.
% ===========================================================================
% splitref.tex -- the within-well split-half reference, as the three splits are
% actually scored and with the resolved +0.1086 reliability line imposed on all
% three.
% \input by Sections/05-Limitations.tex. Compiles against thesis_v2/preamble.tex
% and nothing else. Every number in the drawing is written by
% figs/splitref_numbers.py into figs/fig-splitref.json; run that script if the
% source records ever change, and paste its "TikZ dot values" block below.
%
% STALE GENERATOR -- READ BEFORE RE-RUNNING. splitref_numbers.py:57 hardcodes
% REPRO_LINE = 0.109 and its own comment (:54-56) says 0.109 was chosen because
% it "reproduces all three printed values". That is backwards: the line was
% fitted to the prose rather than taken from the definition, and the prose value
% it was fitted to is the one that is wrong. Panel B below is now drawn at the
% resolved +0.1086 line the document actually defines (see THE LINE). Re-running
% the script as it stands will silently revert this file. Set REPRO_LINE
% = 0.1086 there first, and regenerate figs/fig-splitref.json, which still
% records the superseded 0.109 reconstruction.
%
% v1's version sat inline in Sections/05-Limitations.tex at lines 193-249.
% Four things were wrong with it and are fixed here.
%
%   1. IT PAID FOR FULLWIDTH AND DID NOT USE IT. The old float opened
%      \begin{fullwidth} -- 174mm, and by the preamble's own rule the page
%      carrying one may then carry no margin note -- and drew 150.2mm. It shows
%      six numbers, all six of which its own caption already states, so it was
%      buying 32mm of extra measure to repeat the caption. This draws inside the
%      142mm body measure (see MEASURED below) and that page keeps its margin
%      column.
%
%   2. THE RIGHT-HAND BRACE WAS A DEGENERATE GLYPH. Three values spanning
%      0.956-0.961 on an axis running 0.82-0.98 collapse a
%      decorations.pathreplacing brace to about 1.5pt of path, shorter than its
%      own 3.5pt amplitude, and pgf draws the two halves overlapping: it printed
%      as a stray chevron beside "0.006". A brace that renders as a chevron is
%      worse than no brace. Both spans are now square brackets -- a form that
%      degrades gracefully, so the SAME mark is used in both panels and the
%      reader compares 47.8mm of bracket against 1.8mm of bracket directly.
%      Correspondingly the three converging leaders on the right panel's three
%      near-coincident dots are gone: a single block set hard against the
%      cluster lists all three values in the dots' own top-to-bottom order, so
%      it reads as the cluster magnified rather than as a legend.
%
%   3. NEITHER PANEL SAID WHAT WAS ON THE VERTICAL AXIS. The caption called it
%      "the reference" and the drawing named nothing. There is now a rotated
%      axis title naming the estimator and its frame.
%
%   4. THE TRUNCATION WAS SILENT. The axis runs 0.82-0.98 -- 16 points of a
%      100-point scale -- and that truncation is most of why the 0.132 spread
%      looks dramatic. Both panels share the scale, so the before/after
%      comparison is honest, but the reader has to be told. The axis is drawn
%      BROKEN, with 0.50 (chance, the bottom of the NIR scale's meaningful
%      range) marked below the break, and the caption states the range in words.
%
% GEOMETRY. x and y are both 1mm, so every coordinate below is a millimetre and
% the numbers can be checked against the MEASURED reading. \yv maps a reference
% value onto the axis: \lo..\hi spans \H mm.
%
%   spine            x = 7.5, y = 0 (0.82) .. 58 (0.98)
%   break            gap y = -4 .. 0, two slashes across it, stub -8 .. -4
%   panel A          bracket x = 13, dots x = 36, names from 38.5, values to 66
%   divider          x = 68.5
%   panel B          bracket x = 88, dots x = 91.5, value block from x = 94.5
%   gridlines        stop at x = 94, clear of the value block, uniform length
%
% The two brackets sit on OPPOSITE sides of their dot columns and that is
% deliberate. Panel A's is 47.8mm tall and reads as a measure from anywhere;
% panel B's is 1.8mm and would read as a stray mark if it were the 23mm from
% its dots that panel A's is, so it is set hard against them and its label is
% thrown outward instead. Same mark, same weight, same colour; only the side
% differs, and the two lengths stay directly comparable across the divider.
%
% The two lower dots of panel A sit 0.0127 apart, which is 4.62mm here, against
% a \scriptsize line box of 3.35mm -- so both labels sit at their own dot's
% height and NO leader is needed on the left panel either. That clearance is
% why \H is 58 and not the 52 the old version used.
%
% MEASURED, not estimated. figs/_harness_splitref.tex \input s this file against
% the real preamble, captures the float into a box with \caption and \label
% gobbled, and \typeout s \wd. Reading:
%   tikzpicture  377.33pt = 132.90mm wide, 220.73pt = 77.74mm tall
%   \textwidth   404.03pt = 142.00mm
%   fills 93.4% of the measure, 26.70pt = 9.39mm of slack.
% The build's overfull-hbox gate is 0; this figure contributes none. The
% fullwidth version it replaces drew 150.2mm against a 174mm allowance.
%
% NUMBERS. Panel A is read straight from CANONICAL_NUMBERS.json -> ceilings{}.
% Panel B is reconstructed from re_v3.json -> records[] at repro_cos > 0.1086;
% no file in RESULTS_cluster holds those three values. figs/fig-splitref.json
% carries the full provenance, and is STALE: it still describes the superseded
% 0.109 reconstruction (299 / 394 / 111 kept, train 0.955).
%
% THE LINE. 04c-reencoding.tex:212-213 defines it as the measured null's 95th
% percentile, "+0.109, resolving to +0.1086", and the document is consistent at
% 0.1086, not at 0.109:
%   - 04c:220-221  "Of the 300 scored training conditions ... none falls at or
%                  below the resolved +0.1086." The train minimum repro_cos in
%                  re_v3.json is 0.10892294 -- above 0.1086, BELOW 0.109. At
%                  > 0.109 the figure dropped exactly the condition the text
%                  says is not dropped.
%   - 04c:959      "499 conditions sit below the +0.1086 line" on unseen_combo,
%                  i.e. 396 of 895 kept. The source count of records under
%                  0.1086 is exactly 499; under 0.109 it is 501, giving 394.
%   - 04c:957-958  "train is by construction the reliability-surviving
%                  population", so imposing the line on train must not move
%                  train. At 0.1086 it does not: panel B's train is the same
%                  0.9555006 as panel A's, to the last digit. At 0.109 it moved.
% Recomputed from re_v3.json at repro_cos > 0.1086: kept 300 / 396 / 111, means
% 0.9555006 (train), 0.9605642 (unseen_combo), 0.9597734 (unseen_drug).
%
% TWO NUMBERS THIS FIGURE PRINTS THAT THE PROSE CONTRADICTS. Neither is
% arbitrated here; both are reported upward. The figure does not invent a third
% reading of either.
%   (a) train, panel B. Drawn and printed 0.956, the value the resolved line
%       gives. The prose still prints 0.955 at 04c-reencoding.tex:883 and
%       05-Limitations.tex:224, both inherited from the 0.109 reconstruction.
%       Those two sites need the same correction; they are outside this file.
%   (b) the right-hand spread, printed 0.006. Under the old 0.109 line that was
%       max(rounded) - min(rounded) = 0.961 - 0.955, the raw spread being
%       0.00517. Under the resolved line BOTH conventions give 0.005: the raw
%       spread is 0.0050636, and max(rounded) - min(rounded) = 0.961 - 0.956.
%       So 0.006 is no longer derivable at the line the document defines. It is
%       nevertheless what the prose prints in three places
%       (04c-reencoding.tex:884, 05-Limitations.tex:225 and 03-Apparatus.tex:378,
%       the last a table note in the apparatus chapter), so it is left standing
%       here rather than silently changed to 0.005 -- changing it would put the
%       figure at odds with the prose in three places, and the spread is a
%       claim. FLAGGED, NOT FIXED.
% The left-hand 0.132 is the raw spread rounded and is unaffected.
% ===========================================================================
\begin{figure}[tbp]
% V7: keep the long caption in the float, below the plot, so its complete
% height participates in page placement and cannot run off the margin.
\begin{fullwidth}
\centering
\captionsetup{margin=0pt}
\begin{adjustbox}{max width=\linewidth}
\begin{tikzpicture}[x=1mm, y=1mm, font=\small]
  % --- the constants every coordinate is derived from ----------------------
  \def\lo{0.82}\def\hi{0.98}\def\H{58}
  \newcommand{\yv}[1]{{(#1-\lo)/(\hi-\lo)*\H}}
  \def\ax{7.5}          % spine
  \def\bra{13}\def\dota{36}\def\laba{38.5}\def\vala{66}
  \def\divx{68.5}
  \def\brb{88}\def\dotb{91.5}\def\blk{94.5}
  \def\srf{2}           % bracket serif length, both panels

  \tikzset{
    grid/.style   ={draw=rulegrey!45, line width=0.3pt},
    spine/.style  ={draw=rulegrey, line width=0.5pt},
    tick/.style   ={font=\scriptsize, text=quietink},
    dot/.style    ={fill=ink},
    % text height/depth normalised so a \texttt name and a maths value set at
    % the same y sit on the same baseline; without it Inconsolata's box and the
    % maths box differ by 0.7pt and the three rows print visibly askew.
    lbl/.style    ={font=\scriptsize, text=ink,
                    text height=1.55ex, text depth=0.30ex},
    blk/.style    ={font=\scriptsize, text=ink},
    span/.style   ={draw=accent, line width=0.5pt},
    strong/.style ={font=\scriptsize\bfseries, text=accent}}

  % ---- 1. the shared scale ------------------------------------------------
  % One axis serves both panels; that is the argument, so it is drawn once.
  \foreach \v in {0.82,0.86,0.90,0.94,0.98}{
    \draw[grid] (\ax,\yv{\v}) -- (94,\yv{\v});
    \node[tick, anchor=east] at (\ax-1.5,\yv{\v}) {$\v$};}

  \draw[spine] (\ax,0) -- (\ax,\H);
  % the break, and chance below it. 0.50 is the foot of the NIR scale's
  % meaningful range and it is 32 points below anything drawn, so it is marked
  % past a break rather than pretended onto the plotted span.
  \draw[spine] (\ax,-8) -- (\ax,-4);
  \draw[spine] (\ax-1.3,-3.0) -- (\ax+1.3,-2.2);
  \draw[spine] (\ax-1.3,-1.8) -- (\ax+1.3,-1.0);
  \node[tick, anchor=east] at (\ax-1.5,-8) {$0.50$};
  \node[tick, anchor=west] at (\ax+1.5,-8) {chance};

  \node[rotate=90, anchor=south, font=\scriptsize, text=quietink, align=center]
    at (-1,29) {within-well split-half reference\\(residual frame, $\NIR$)};

  \draw[draw=rulegrey!60, line width=0.4pt] (\divx,-1) -- (\divx,60);

  \node[strong, anchor=south] at (39.5,61.5) {as the splits are actually scored};
  \node[strong, anchor=south] at (94,61.5)
    {with the training line imposed on all three};

  % ---- 2. panel A: the reference as each split is actually scored ---------
  % values: CANONICAL_NUMBERS.json -> ceilings{}
  % Name anchored west at \laba, value anchored EAST at \vala, so the three
  % values line up in a column despite the three split names differing by six
  % characters. Left-setting the whole label leaves them ragged.
  \fill[dot] (\dota,\yv{0.955501}) circle (1.4pt);
  \fill[dot] (\dota,\yv{0.836477}) circle (1.4pt);
  \fill[dot] (\dota,\yv{0.823720}) circle (1.4pt);
  \node[lbl, anchor=west] at (\laba,\yv{0.955501}) {\texttt{train}};
  \node[lbl, anchor=west] at (\laba,\yv{0.836477}) {\texttt{unseen\_drug}};
  \node[lbl, anchor=west] at (\laba,\yv{0.823720}) {\texttt{unseen\_combo}};
  \node[lbl, anchor=east] at (\vala,\yv{0.955501}) {$0.956$};
  \node[lbl, anchor=east] at (\vala,\yv{0.836477}) {$0.836$};
  \node[lbl, anchor=east] at (\vala,\yv{0.823720}) {$0.824$};

  \draw[span] (\bra+\srf,\yv{0.955501}) -- (\bra,\yv{0.955501})
           -- (\bra,\yv{0.823720}) -- (\bra+\srf,\yv{0.823720});
  % 0.889610 and 0.958032 below are the MIDPOINTS of the two spans, used to
  % place each label at the middle of its own bracket. They are geometry, not
  % data, and appear nowhere in fig-splitref.json.
  \node[strong, anchor=west] at (\bra+\srf+1.5,\yv{0.889610}) {spread $0.132$};

  % ---- 3. panel B: the same reference, training line imposed on all three --
  % values: reconstructed from re_v3.json -> records[] at repro_cos > 0.1086,
  % the resolved line of 04c-reencoding.tex:212-213. Kept 300 / 396 / 111.
  % train is 0.955501 -- byte-identical to panel A's train above, because all
  % 300 train conditions clear the line. That identity is the point: the line
  % selects train, so imposing it on train cannot move train.
  \fill[dot] (\dotb,\yv{0.960564}) circle (1.4pt);
  \fill[dot] (\dotb,\yv{0.959773}) circle (1.4pt);
  \fill[dot] (\dotb,\yv{0.955501}) circle (1.4pt);

  % Bracket hard against the dots, opening toward them; label thrown left.
  \draw[span] (\brb+\srf,\yv{0.960564}) -- (\brb,\yv{0.960564})
           -- (\brb,\yv{0.955501}) -- (\brb+\srf,\yv{0.955501});
  % 0.006 is the prose's number, not this panel's: see (b) in the header. The
  % bracket it labels spans 0.0050636, i.e. 0.005.
  \node[strong, anchor=east] at (\brb-1.5,\yv{0.958032}) {spread $0.005$};

  % Three dots 1.8mm apart cannot carry three leaders, and a leader short enough
  % to reach a block set beside them is indistinguishable from a hyphen. So: no
  % leader. One block, hard against the cluster with nothing else within 25mm of
  % it, listing all three in the dots' own top-to-bottom order, so it reads as
  % the cluster magnified rather than as a legend.
  \node[blk, anchor=west] at (\blk,\yv{0.958032})
    {\begin{tabular}{@{}l@{\hspace{2mm}}l@{}}
       \texttt{unseen\_combo} & $0.961$\\
       \texttt{unseen\_drug}  & $0.960$\\
       \texttt{train}         & $0.956$
     \end{tabular}};
\end{tikzpicture}%
\end{adjustbox}
\caption[The split references, before and after matching]{\captiontitle{The three
splits do not differ in recoverable signal; they differ in which conditions were
allowed into them.} The vertical axis is the within-well split-half precision
reference --- one half of a condition's cells scored against the other half's
truth, in the residual frame, on the $\NIR$ scale --- and it is
\emph{truncated} to $0.82$--$0.98$, a sixth of the scale; chance is $0.50$,
marked below the axis break and far under everything drawn. Both panels share
that truncated scale, so the two spans may be read against each other. Left,
the reference as each split is actually scored, $0.956$ on \texttt{train}
against $0.824$ and $0.836$ on the held-out splits, a spread of $0.132$. Right,
the same reference with the resolved $+0.1086$ reliability line imposed on all
three, $0.956$, $0.961$ and $0.960$, a spread of $0.005$; at this scale the
three sit inside a bracket two millimetres tall, so they are listed beside the
cluster in
their own order rather than labelled in place. The left-hand spread is a
selection artefact. The markers are bare: no interval is carried for this
reference in the source record, and none is drawn rather than invented.\label{fig:splitref}}
\end{fullwidth}
\end{figure}


\emph{Those conditions carry less signal} is therefore not a reason to discount
a held-out result, and this thesis nowhere uses it as one. What survives is
narrower and binding: the matched-population rescore covered the precision
reference only. The model's own arms were never rescored, so where its held-out
scores would sit once the populations agree is not reported here, and the
train-against-held-out comparison of \cref{sec:res-generalisation} inherits the
mismatch --- which matters most for the score that section reports as covering
chance on both held-out splits.

% PLACEMENT ONLY. Eight lines and a framed header, and it was splitting five and
% three across the page turn, so the header sat with half its own text and the
% rest arrived under the running head as an unlabelled fragment. 11 lines is the
% box entire.
\needspace{11\baselineskip}%
\begin{scopelimit}[discharges the limit carried from \cref{sec:res-lookup}]
Three selection rules are in force at once, and no result computed under one of
them may be set beside a result computed under another. A raw score read across
the three splits is not answering the same question on each, because the splits
were filtered by different rules. An absolute $\NIR$ computed with the
comparison set grouped by cell line is not comparable with one computed within
plate. And a per-drug lookup is the bar for a drug that was \emph{seen} in
training and is being asked about a new context: no lookup exists for a drug
that never appeared in training, so that regime keeps its own bar, set in
\cref{sec:res-scope}.
\end{scopelimit}

% ===========================================================================
\section{What the intervals do not cover}
\label{sec:lim-intervals}

\paragraph{No training run was repeated} Seeds in this work move prompts,
partners and resampling. None of them moves the weights. Every arm here was
trained once, from one base checkpoint under identical hyperparameters
(\cref{tab:trainconfig}), so each dissociation this thesis rests on --- the
residual re-encoding against the standard target, and the optimal-transport
target against both --- compares one training run with one training run. A
second run of the same arm under a different data order would land on a
different checkpoint and score differently, and the size of that difference was
never measured here. It is therefore not covered by any interval in this
document: the intervals quantify sampling over conditions, cell lines and wells,
and a between-arm gap of the size reported would survive or fail a training
replication that nobody has run. The ten-epoch diagnostic varies the horizon and
not the seed, so it bounds undertraining and not run-to-run spread. What this
costs is specific rather than general: the \emph{direction} of the re-encoding
effect is corroborated by the token-divergence measurement, which needs no
model at all, but the \emph{magnitude} of every arm-to-arm difference is an
n-of-one comparison.

\paragraph{A swing that looked like generation variance and was not} Unseeded
generation is the kind of defect that surfaces as an unstable number, and once
it seemed to.

\begin{withdrawn}{seed instability in the held-out-drug control}
Two superseded evaluations of the held-out-drug control returned values of
$\gap$ far enough apart to be recorded as evidence of seed instability. Only one
of the pair survives in the results archive, so the size of the swing is no
longer quotable. The diagnosis was wrong. It named the trigger and missed the
cause.
\end{withdrawn}

Both of the superseded figures are gaps against the \texttt{opposite} swap, which
\cref{sec:res-generalisation} shows is not a null; it carries a term for how far
the \emph{partner} drug was pushed the wrong way, which changes with the
draw.\seealso{\cref{sec:res-generalisation}, where a neutral stratum replaces
\texttt{opposite} as the comparator} On the canonical evaluation the same
control reads $+0.0379$ against \texttt{opposite} and
\ci{-0.0315}{-0.0836}{+0.0206} against the neutral stratum, clustered two-way on
cell line and well.

Two mechanisms were at work, and both are now addressed. \texttt{do\_sample}
advanced from global torch state, so the two arms of a contrast were drawn from
different points of one
stream; each call now seeds its own generator from (condition, arm, batch), and
the comparator has been corrected. The general caution stands regardless; this
swing is no longer evidence for it.

\paragraph{The intervals omit generation variance} The intervals here are
cluster-robust with each condition's score treated as fixed: two-way over
\emph{cell line} and \emph{treated well} on a $\min(G_{\textrm{line}},
G_{\textrm{well}}) - 1$ Student $t$ reference for the headline gaps, reported
again clustered on the \emph{drug}; multiway dyadic for the transfer
coefficient; a one-way clustered bootstrap, unit named, in a few tables.

They capture between-cluster variance and leave out the variance from sampling
the model's predictions. With four sampled generations per condition at temperature
$0.8$ in the residual-frame arms and eight in the expression-frame ones
(\cref{sec:methods-data}), that source is not negligible, and no across-seed
replication of the generation arms was run, deferred or planned. Every interval
quoted from a generation arm covers the sampling of conditions, cell lines and
wells, and stops there. That includes the within-condition slope regressions of
\cref{sec:res-blind}, \cref{sec:res-epsilon} and \cref{sec:res-generalisation},
whose points are themselves sampled generations, so the comparator-free
instrument this thesis leans on hardest is bounded the same way as the rest.

One measurement \emph{was} repeated across seeds, recorded so it is not mistaken
for the replication that was not run. The activation probe of \cref{sec:res-mechanism} returned \ci{+0.0197}{+0.0063}{+0.0335} on the anchor seed $42$ and \ci{+0.0118}{+0.0060}{+0.0173} and \ci{+0.0085}{+0.0029}{+0.0143} on seeds $43$ and $45$, each excluding zero. Its spread does not transfer to $\NIR$ and does not
establish a scale for generation variance. Those runs share a checkpoint and
score teacher-forced log-probabilities, so their seed moves the prompts, the
partners and the bootstrap, and leaves both the weights and the sampled
prediction untouched. They are scored on $60$, $240$ and $240$ prompts in that order, so
the largest of the three effects is also the smallest-$n$ one. Each of those three is a single-test interval that carries no correction for the twenty-six tests of \cref{fig:family}, and the first is the draw whose survival that correction leaves in doubt (\cref{sec:res-mechanism}); three uncorrected seeds excluding zero therefore do not reinstate the effect.

\paragraph{Resolution and cross-plate absolutes} Residuals are condition-level
pseudobulks over at least $40$ treated cells, so single-cell heterogeneity of the
response is not modelled in this frame at all; \cref{sec:res-blind} measures what
a single cell can and cannot deliver. Those pseudobulks are then scored against
comparison sets grouped by cell line rather than by plate. The \emph{difference}
$\gap$ is leak-immune, since both arms share the same control cell, but the
absolute $\NIR$ values are a different matter. The plate-scoped generic removes
plate structure in expectation only, and it is estimated from a small plate
group, so some may survive in an individual residual. The rescore that would size
what survives, regrouping the generation-free arms within plate, needs no GPU
time and was not run.

What exists instead is one-sided, and \cref{tab:plateleak} bounds it from
above.\framecue{\Eframe}{the plate-leak audit} In the expression frame with no
generic subtracted, holding the plate constant moves a drug-agnostic mean
baseline from $0.399$ to $0.161$ and the precision reference from $0.625$ to
$0.575$: the yardstick falls with the baseline, so the leak is not an offset one
side of a comparison could be corrected by. Residual-frame absolutes are read
within their own frame and never set beside a within-plate number. The audit's within-plate figures sit slightly off the benchmark figures of \cref{sec:res-blind}: control-copy $0.510$ against $0.504$, mean baseline $0.161$ against $0.180$, precision reference $0.575$ against $0.576$.

\begin{table}[htbp]
% MARGINALIA: in-column at natural width; caption and both notes in the margin.
\begin{lrbox}{\tvfb}
\begin{tabular}{@{}l r r@{}}
\thead{Predictor (drug-agnostic)} & \thead{cross-plate}
  & \thead{within plate} \\
\midrule
mean baseline & 0.399 & \key{0.161} \\
control-copy & 0.551 & 0.510 \\
\addlinespace[10pt]
\quiet{precision reference (within-well split-half)} & 0.625 & 0.575 \\
fraction of drugs at chance & 43\% & 49\% \\
\end{tabular}
\end{lrbox}
\sidecap{table}{\captiontitle{Expression-frame $\NIR$, the leak audit.} A drug-agnostic mean
baseline more than doubles its score when the plate is free to vary. The
precision reference falls once the plate is held constant, because batch
structure was inflating it too.\label{tab:plateleak}
  \par\addvspace{6pt}A wider diagnostic run than the $n = 606$ tier-2 benchmark the blindness
verdict is stated on, at $n = \num{1163}$ cross-plate and $n = 998$ within
plate. Holding the plate constant costs every condition with no same-plate
comparison set, so the within-plate column is a subset of the cross-plate one
rather than the same conditions rescored.}
\end{table}

% "The similarity is a choice too". States what the frame's choice of SIMILARITY
% costs. (Its companion paragraph on the choice of scored OBJECT was removed as
% duplicating sec:res-transfer.) Sources, all read directly:
%   nir_benchmark.py:13     "expr-NIR : Euclidean distance after decoding ranks
%                            -> expression via linear_model.json"
%   nir_benchmark.py:363-4  d_own = np.linalg.norm(pe - truth_expr[d]) etc.
%   nir_benchmark.py:864    "expression_distance": "Euclidean"
%   residual_eval.py:13     "Similarity = cosine."
%   residual_eval.py:125    def cos(a, b)
% The rescore: RESULTS_cluster/nir_sameplate_profiles.npz, 606 predictions and
% truths over 80 groups, tier2_unseen_drugs, the same support as the frozen
% manifest, scored under the definition of sec:res-metrics with the half-tie
% rule. -||.||_2 gives mean 0.4977 median 0.5000; cosine gives mean 0.5007
% median 0.5000; Spearman 0.6217 between the two per condition; 352/606 move by
% more than 0.10. The Euclidean value reproduces the published 0.498 on the same
% n, which is the only thing validating the rescore, so it is stated.
% The Spearman was re-derived from the same npz under the same half-tie rule
% (scipy spearmanr on the 606 per-condition pairs; Pearson 0.6218, Kendall
% 0.4785 on the same pairs). Restricting to groups of at least three drugs and
% dropping the half-tie rule both return 0.6217; averaging Spearman within
% groups returns 0.5775, a different estimator and not the one quoted.
% NOT RESCORABLE: every local .npz is expression-frame. There is no
% re_*_profiles.npz, so the residual frame carries no such measurement and the
% last sentence says so, which is what stops a reader generalising.
% The sigma instances belong at the definition in 04a-ruler.tex, which already
% parameterises by sigma and already declares the expression instance; this pass
% was scoped to one file, so the residual instance is named here instead.
\paragraph{The similarity is a choice too} $\NIR$ is parameterised by a
similarity $\sigma$ (\cref{sec:res-metrics}), and the two frames instantiate it
differently: the expression frame scores $\sigma(\hat y, y) = -\lVert \hat y - y
\rVert$, the residual frame the cosine. What that choice carries was measured on
the run the blindness verdict is stated on. The same $606$ stored predictions
and truths, the same $80$ comparison groups, tier 2 within plate, were rescored
under both similarities using the definition above, half-tie rule included: the
aggregate reads $0.4977$ under the negated Euclidean distance against $0.5007$
under the cosine, both at chance. The Euclidean value reproduces the published
figure on that support, which is what validates the rescore.

What does not hold across the substitution is the ordering beneath the
aggregate. The two scorings agree per condition at Spearman $0.6217$, and $352$
of the $606$ conditions move by more than $0.10$. That neither strengthens nor
weakens the aggregate verdict: neither similarity establishes aggregate
discrimination on this run. The per-condition changes need not be noise;
an aggregate near chance can conceal subgroup gains and failures, and the
rescore does not establish which conditions carry either.

Two properties of the published convention belong beside it. A Euclidean
distance is dominated by its largest coordinates, so the high-abundance genes
carry the comparison. And the vectors compared are reconstructions: the
prediction is a cell sentence, decoded to expression through a fitted linear map
before any distance is taken, and that map was fitted for absolute fidelity and
never validated for discrimination (\cref{sec:lit-c2s}). The residual frame
carries no such measurement, because its predictions and truths were not stored
in a form that can be rescored, so nothing is claimed about its sensitivity to
$\sigma$ in either direction.

\begin{scopelimit}[carried into \cref{sec:res-blind} and every per-condition
reading of an expression-frame score]
An at-chance aggregate $\NIR$ in the expression frame holds under either
similarity, negated Euclidean or cosine, on the one run rescored here, tier 2
within plate. Expression-frame aggregates that carry signal were not rescored,
so their standing under $\sigma$ is untested. No per-condition or per-stratum
reading has that standing either: on the run that carries the verdict the two
orderings agree only at Spearman $0.6217$, so no figure below the aggregate may
be presented as independent of $\sigma$. The residual frame carries no such
measurement, so its scores are covered neither way.
\end{scopelimit}

\section{The unseen-drug arm}
\label{sec:lim-unseen}

Two limits on the unseen-drug arm must be stated: what the prompt reveals and
what the intervals resolve. A separate defect, the comparator, was found and
corrected, and \cref{sec:res-generalisation} reports the arm against the neutral
stratum throughout.

\paragraph{The design leaks the drug's class} Holding out a drug does not
remove everything the model knows about it. The prompt contains
\texttt{Mechanism: \{moa\}}, so a held-out drug still arrives with its class
label. The correct design is leave-one-\emph{mechanism}-out.

The leak is measured, and the measurement is not clean. Scrambling the
\texttt{Mechanism} field alone moves the model's $\NIR$ by
\ci{+0.0303}{-0.0078}{+0.0685} ($n = \num{738}$, clustered on drug and well), an
interval containing zero. That understates the channel; it does not bound it,
because this row's swap partner scores $0.520$, above chance, and a comparator
above chance suppresses the gap (\cref{tab:fieldattribution}, Panel~B). A
mechanism-only lookup is the same measurement from the other side and equally
inconclusive. Pooled over the $n = \num{218}$ conditions on which both are
scored, model $-$ \texttt{moa\_lookup} $=$ \ci{+0.044}{-0.048}{+0.135},
clustered two-way on cell line and well, spans zero, so the comparison settles
nothing in either direction.

On \texttt{unseen\_drug}, the split where it would matter, the mechanism lookup
is ahead of the model by $0.099$ on $44$ conditions; in the contrast's reported
direction, model minus lookup is \ci{-0.099}{-0.222}{+0.008}, clustered on cell
line. The interval spans zero and does not establish which predictor is better
on that support. It also does not establish how much the mechanism field
contributes to either score. The arm tests held-out drug identities with
mechanism information still supplied, so it cannot establish transfer to a new
mechanism family.

\paragraph{The gap does not establish equivalence to zero} On the $n = 199$
\texttt{unseen\_drug} conditions, the model-minus-neutral-scramble gap is
\ci{-0.0315}{-0.0836}{+0.0206} clustered two-way on cell line and well, and
\ci{-0.0315}{-0.1049}{+0.0419} clustered on drug (\cref{fig:power}). Both
intervals include zero and a range of small positive and negative effects.
Neither establishes a positive unseen-drug gap, but neither is an equivalence
test or evidence that transfer is absent.

The pooled mechanism-only lookup scores $0.552$, or $0.052$ above chance.
That is a baseline result on different support, not an upper bound on the
effect a mechanism-conditioned model could achieve. It is also a different
contrast from the model-minus-scramble gap; it cannot supply the gap's maximum
plausible effect. Likewise, an approximate power calculation based on the
model's absolute $\NIR$ interval does not determine which effects lie inside
the observed gap intervals.

% ===========================================================================
% Unseen-drug gap: one estimand, two reported 95% confidence intervals.
% Exact values and provenance are recorded in the sibling power.json.
% Geometry: 34mm row labels + 3mm gutter + 86.4mm axis = 123.4mm nominal.
% The pooled lookup and an MDE derived from absolute NIR are not this gap.
% ===========================================================================
\begin{figure}[htbp]
\begin{fullwidth}\raggedright
\begin{adjustbox}{max width=\linewidth}
\begin{tikzpicture}[
    x=480mm, y=1mm, font=\footnotesize,
    rowlbl/.style={text=quietink, align=flush right, text width=34mm,
                   inner sep=0pt, anchor=east, xshift=-3mm},
    ticklbl/.style={font=\scriptsize, text=quietink, inner sep=0pt,
                    anchor=north},
    note/.style={text=quietink, inner sep=0pt}]

  \draw[rulegrey, line width=0.5pt] (-0.12,0) -- (0.06,0);
  \foreach \t/\lab in {%
      -0.10/$-0.10$, -0.05/$-0.05$, 0/$0$, 0.05/$+0.05$}{%
    \draw[rulegrey, line width=0.5pt] (\t,0) -- (\t,-1.4);
    \node[ticklbl] at (\t,-2.1) {\lab};}
  \node[note, font=\scriptsize, anchor=north, align=center]
    at (-0.03,-7.5)
    {model $\NIR$ minus neutral-scramble $\NIR$\\
     ($0$ = no difference)};

  % Each zero segment crosses its own interval, clear of the numeric label.
  \draw[rulegrey, dashed, line width=0.5pt] (0,28.5) -- (0,37);
  \draw[rulegrey, dashed, line width=0.5pt] (0,10.5) -- (0,19);
  \node[note, anchor=south, text=ink] at (-0.03,42)
    {estimate $-0.0315$\enspace;\enspace$n = 199$};

  % Two-way clustering on cell line and treated well.
  \node[rowlbl] at (-0.12,32) {cell line and well\\(two-way)};
  \draw[ink, line width=0.9pt]
    (-0.08358077892636565,32) -- (0.020617264814239943,32);
  \draw[ink, line width=0.8pt]
    (-0.08358077892636565,30) -- (-0.08358077892636565,34);
  \draw[ink, line width=0.8pt]
    (0.020617264814239943,30) -- (0.020617264814239943,34);
  \fill[accent] (-0.03148175705606285,32) circle (1.7pt);
  \node[note, anchor=north west] at (-0.12,26.5)
    {95\% CI $[-0.0836,\,+0.0206]$};

  % Clustering on drug; the point estimate and scored conditions are unchanged.
  \node[rowlbl] at (-0.12,14) {drug};
  \draw[ink, line width=0.9pt]
    (-0.10487258422550776,14) -- (0.04190907011338206,14);
  \draw[ink, line width=0.8pt]
    (-0.10487258422550776,12) -- (-0.10487258422550776,16);
  \draw[ink, line width=0.8pt]
    (0.04190907011338206,12) -- (0.04190907011338206,16);
  \fill[accent] (-0.03148175705606285,14) circle (1.7pt);
  \node[note, anchor=north west] at (-0.12,8.5)
    {95\% CI $[-0.1049,\,+0.0419]$};
\end{tikzpicture}
\end{adjustbox}
\caption[The unseen-drug gap under two clusterings]{\captiontitle{The
unseen-drug gap includes zero under both clusterings.} The points show the same
model-minus-neutral-scramble estimate on $n = 199$ conditions in the
\texttt{unseen\_drug} split, $-0.0315$. The comparator is the
geometry-defined \texttt{orth} scramble on those conditions. The upper row
shows the 95\% interval clustered two-way on cell line and treated well,
$[-0.0836,\,+0.0206]$; the lower row clusters on drug,
$[-0.1049,\,+0.0419]$. Both intervals include zero and a range of small
effects. Neither establishes a positive gap or equivalence to zero. They
quantify the reported cluster uncertainty, conditional on the trained
checkpoint and sampled generations (\cref{sec:lim-intervals}).}
\label{fig:power}
\end{fullwidth}
\end{figure}


As an \emph{instrument control}, this result is compatible with the neutral
comparison in \cref{sec:res-generalisation}; compatibility with zero alone does
not validate a comparator or confirm a null effect. ``Unseen'' is narrow in a
second way too. Pythia-1b's pretraining is natural-language text in which drug
names occur, and the prompt still supplies \texttt{Mechanism}.

\begin{scopelimit}[carried into \cref{sec:conclusions}]
The \texttt{unseen\_drug} arm is read as an instrument control. Its estimated
gap is compatible with zero and a range of small effects under both reported
clusterings. This does not establish equivalence to zero or the absence of
unseen-drug transfer. ``Unseen'' means held out from Tahoe fine-tuning; it does
not mean a new mechanism family or absence from natural-language pretraining.
\end{scopelimit}

% ===========================================================================
\section{Future work}
\label{sec:lim-future}

A limitation earns its place when something could be run to remove it.

\paragraph{The arm to build carefully} The obvious continuation of
\cref{sec:res-scope} is a channel-conditioned prompt that appends protein targets
to the instruction, retrains, and evaluates on a leave-one-mechanism-out split.
We do not recommend building it eagerly. \Cref{sec:res-field} leaves unresolved
which span of the prompt is even read, so a new field would be added to an
instruction whose existing fields are of unmeasured value; and the case against
it does not rest on that anyway, but on the size of the readout deficit. On the
common support the model covers $13.9\%$ of the distance from chance to the
within-well bar that a context-blind lookup covers entirely, and adding a field
to a prompt does not touch that deficit.

Worth building instead is an intervention that makes the drug \emph{causal}
rather than merely present: a learned per-drug modulation of the transformation
from control to treated, so the drug parameterises the computation instead of
sitting in the residual stream as an input the readout can round away. That is a
research architecture, outside the scope of what remained here.

\paragraph{Two channel questions that remain open} Target and mechanism overlap,
so whether target adds anything \emph{over} mechanism is unresolved, and it
matters: the prompt already exposes mechanism, and whether the model reads that
field is itself unsettled. Second, the gate of \cref{sec:res-scope}
measures \emph{similarity transfer}, an unseen drug resembling a seen drug with
the same target. It does not test a \emph{compositional rule}, that a named
target's own pathway moves, which would apply to any drug whose target is named,
including one with no similar drug in training. That claim is the stronger one,
it does not reduce to a lookup, and it is untested.

Protein target clears the relevance margin on the cell-line axis, misses it on
the drug axis against both estimand-correct nulls, clearing it there only
against the weaker count-matched one, and rests on $18$ annotated drugs of which
two are held out; mechanism misses its lower bound by $0.0021$ on the cell-line
axis and spans zero on the drug axis; chemistry's interval straddles the margin
on both. Chemistry attenuates under the different-plate construction, but that
construction also changes support and partner availability, so the attenuation
cannot be assigned causally to co-plating. On its restricted support, same-target
retrieval is the strongest observed lead, and stands only as a direction worth
testing. The binding limitation is coverage: $18$ drugs cannot settle it, and
the experiment that would is a target-annotated panel broad enough to cluster
on.

\paragraph{Unseen cell lines} A cell line never seen with any drug is a
different generalisation axis from the held-out wells tested here, and one the
plating design does not compromise in the same way. It asks whether the control
cell's state is used \emph{compositionally} with the
drug.\seealso{\cref{sec:res-transfer}: $\Tcoef$, and what its complement does
not measure}

The transfer coefficient scopes it only loosely. \ci{44.7\%}{40.7\%}{48.6\%} of
the similarity scale $\Tcoef$ is measured on is not shared across the compared
contexts, and no structure was recovered in that unshared part by tests that do
not settle it. That figure is not a drug$\times$cell-line interaction share,
since it mixes cell line, dose, well and estimator (\cref{sec:res-transfer}), so
it bounds what a per-drug constant leaves on the table without establishing that
a new cell line is the axis the remainder lies along.

No consistent per-line direction was found, so the achievable target on a new
cell line plausibly reduces to recalling the drug's per-drug constant, which a
training-only lookup already does at $0.890$ on held-out wells, the nearest
measured analogue of a new context ($0.913$ pooled, inflated on trained
conditions by the well advantage). The axis is less promising than it appeared,
though the argument is weaker than a coefficient-based one.

\paragraph{Objectives that were considered and rejected} Three training
objectives were designed and then set aside on their own mechanics, recorded so
they are not re-litigated.
\begin{itemize}
  \item \textbf{Reinforcement learning with a contrastive reward.} Its
        per-sequence reward keeps the token-dilution ratio out of the gradient,
        attacking a bottleneck the residual encoding has already lifted ($34.6
        \rightarrow 118.2$ differing tokens at cell-line scope, $120.3$ at the
        plate scope the trained build uses). Its reward optimises similarity to
        the true residual, which \texttt{drug\_lookup} already achieves at
        $0.913$ pooled, $0.890$ on held-out wells. Its best case is convergence
        to the lookup.
  \item \textbf{A differentiable profile-level contrastive loss.} The natural
        construction, $\textrm{score}(g) = \sum_i P(\textrm{token}_i =
        g)\,w(i)$, presumes gene $\approx$ token, but $86.9\%$ of the $\num{946}$
        panel genes share a first sub-token under this tokeniser (mean $3.254$
        BPE tokens per gene, only $8$ single-token). Under teacher forcing the
        target \emph{is} the drug-specific residual, so the true prefix
        identifies the drug within a few symbols and the gradient vanishes at
        the readout solution the loss exists to improve on.
  \item \textbf{Loss re-weighting toward differentially expressed genes.}
        A re-weighted cross-entropy against a single target contains no
        negatives, so nothing forces one drug's prediction to differ from
        another's. It leaves the discrimination failure where it was, and that
        failure is
        what the measurements establish, whether or not the readout is also
        blind to direction specifically, which \cref{sec:res-mechanism} declines
        to establish.
\end{itemize}

\paragraph{One model size, and the smallest one} Every result here comes from a
1B-parameter backbone. C2S-Scale releases checkpoints up to Gemma-2 at 27B, but
these larger models were not evaluated. This thesis therefore does not establish
whether increasing model size improves counterfactual fidelity.

The constraint this work points to is not capacity. The target tokens are
overwhelmingly shared across drugs --- only $17\%$ of them differ between two
drugs in the same cell line under the standard target, against $59\%$ under the
residual encoding at cell-line scope and $60\%$ at the plate scope the trained
build uses (\cref{tab:divergence}) --- so very little of the training signal
can carry drug identity. That step is a hypothesis about the optimisation; no
token-level loss or gradient attribution was computed, so no fraction of the
gradient is claimed. A larger model receives the same diluted signal from the
same objective, and re-encoding lifted the effect while leaving
capacity untouched, which points away from capacity without establishing what it
points towards. But the re-encoding changed five things at once, aggregation,
control referencing, sign coding and length among them, and a budget-matched
optimal-transport target lifts the effect too, so \cref{sec:res-encoding}
declines to identify the tokenisation as the constraint that binds.

The genuinely different argument for scale is \emph{prior knowledge}. A 27B
model plausibly knows, as a fact about the world, that a named drug inhibits a
named kinase, which a 1B model may lack. \cref{sec:res-field} leaves
unresolved whether the model uses the \emph{mechanism} field it already has, at
\ci{+0.0303}{-0.0078}{+0.0685} ($n = \num{738}$) with a comparator above chance
suppressing the gap, and real pharmacological knowledge is exactly the case
where that might change. Testing it needs only a low-rank adaptation of a
mid-sized checkpoint. We did not run it, and it is the first thing we would run
with more time.

\paragraph{One epoch, and a thin slice of the atlas} Data scale is the other
half of that question and it is not the same one. Fine-tuning saw
$\num{620564}$ cells, $\num{600000}$ treated and $\num{20564}$ control, drawn
from the roughly one hundred million Tahoe-100M holds, covering $106$ of its $379$ drugs, for one epoch at one learning rate
(\cref{tab:trainconfig}). Under a hundredth of the atlas was used, and no arm saw
any cell twice. The diluted-signal argument above survives that: a model trained
on ten times the cells receives the same target strings, so more data of the same
kind does not make the drug easier to see in the loss. The ten-epoch diagnostic separately excludes undertraining as a trivial explanation --- validation loss is best at the first epoch and worse at every epoch after it --- though it tests one schedule at one learning rate, and no long-run checkpoint was scored with $\NIR$ (\cref{sec:app-training}).
Neither argument rules out the possibility that a drug must be seen in more
cellular contexts than this sample provides before its response becomes
learnable, and that possibility is \notestablished either way. It is a cheaper
experiment than a larger backbone, and it is not one we ran.

\paragraph{The starting checkpoint was never prepared for this task} We
fine-tune from \texttt{C2S-Scale-Pythia-1b-pt}, whose cell-sentence pretraining
corpus is atlas data alone, CELLxGENE and the Human Cell Atlas, with no
perturbation corpus of any kind, and whose released capabilities do not include
perturbation prediction.

That is why the held-out splits are as clean as they are. A drug withheld here
entered neither the fine-tuning data nor the generic that defines its target,
the split being assigned from metadata before any expression was aggregated, so
no Tahoe perturbation response for it can have been memorised. But the failure
we report is then a failure of \emph{this} starting point, and the fair
follow-up is whether a checkpoint already tuned for perturbation prediction
would behave differently.

Within the C2S-Scale line that question cannot be answered from what is public.
The perturbation-specialised model there is a Gemma-3 4B checkpoint that was not
released. Outside that line it is answerable, and we did not answer it.
Tahoe-x1~\cite{tahoex1} is a perturbation-trained foundation model on this same
atlas, public before this thesis was submitted. It is not a cell-sentence model,
so it does not isolate the representation variable, but it is the obvious
external check on whether a readout failure of this kind survives
perturbation-specific pretraining, and it was not run here. That is a gap in
this work, not a gap in the public record.

Every arm starts from the same base and differs only in the variable under test,
with one material exception, the consensus arm, stopped at roughly $60\%$ of an
epoch with its tier-1 $\NIR$ unscoreable (\cref{tab:trainconfig}). The
\emph{comparisons} between arms are therefore unaffected even where the absolute
level is not, and that arm's null is quoted with the qualification wherever it
is used.

\paragraph{Generality} With one backbone, atlas and representation, this study
leaves open whether readout-specificity failure is specific to cell-sentence
models or general to autoregressive conditional transcriptome generation.
The protocol in \cref{sec:res-used,sec:res-mechanism} (purified subspace, removal
gate, positive control, matched-norm swap) transfers unchanged to models with
an accessible residual stream.
